\section{Reproducibility}\label{reproducibility}

All code, the dates collected through the GitHub API (already
processed), the study table, the trained models and the result files are
available at \url{https://github.com/podorize/skillreuse}. The data come
from a fixed revision of GitSkills
(\texttt{ebab17454a7c236f8f26b183567f1a126f42e3f8}), verified by
SHA-256. The environment is Python 3.12 with every dependency pinned in
\texttt{uv.lock}. All random choices use fixed seeds: \texttt{20261121}
for the sample, \texttt{20261122} for the split, 2027--2029 for the
cross-validation folds, 2027 for models, searches and bootstrap, and
\texttt{20261123} for the interpretation sample.

Start with the Colab notebook
(\url{https://colab.research.google.com/github/podorize/skillreuse/blob/main/notebooks/reproduce.ipynb}):
it follows this paper section by section, shows the code of each step
and runs it on the published data. Locally, after \texttt{uv\ sync},
there are four levels, one command each (times on a laptop with 8 GB of
memory):

\begin{enumerate}
\def\labelenumi{\arabic{enumi}.}
\tightlist
\item
  \emph{Check the results} (\texttt{make\ verify}, under a minute):
  applies the trained models to the holdout and checks every number in
  this paper.
\item
  \emph{Redo the analysis} (\texttt{make\ rerun}, about 20 minutes):
  repeats every step, from the exploratory analysis to the
  interpretation.
\item
  \emph{Redo the processing} (\texttt{make\ rebuild}, about 30 minutes
  plus a 7 GB download): also rebuilds the sample and the study table
  from GitSkills.
\item
  \emph{Redo from scratch} (\texttt{make\ from-scratch}, about 8 hours):
  also dates the copies again through the GitHub API; the result is
  close but not identical, as repositories have changed.
\end{enumerate}
